% Journal manuscript draft -- Physical Review Accelerators and Beams (REVTeX 4.2)
%
% Build:   cd paper/manuscript && latexmk -pdf plasma_column_prab.tex
%          (or pdflatex; bibtex; pdflatex; pdflatex)
% Numbers: every \num... macro is generated by
%          python scripts/make_manuscript_figures.py
%          Do not type physics numbers into this file by hand.
%
% STATUS: DRAFT. Sections marked \pending{} require production PIC results
% that do not yet exist in this repository. See Sec. V and the
% "Status of this draft" note at the end of the document.

\documentclass[aps,prab,10pt,twocolumn,superscriptaddress,showpacs,floatfix]{revtex4-2}

\usepackage{amsmath,amssymb}
\usepackage{graphicx}
\usepackage{booktabs}
\usepackage{xcolor}
\usepackage[colorlinks=true,linkcolor=blue!60!black,citecolor=blue!60!black,urlcolor=blue!60!black]{hyperref}

\graphicspath{{figures/}}
% Auto-generated by scripts/make_manuscript_figures.py -- do not edit by hand.
% Values are analytical design estimates from plasma_column.design_estimates.
\newcommand{\numbeta}{\ensuremath{0.00800}}
\newcommand{\numvbeam}{\ensuremath{2.40\times10^{6}}}
\newcommand{\numKzero}{\ensuremath{1.25\times10^{-3}}}
\newcommand{\numKzeropeak}{\ensuremath{6.25\times10^{-3}}}
\newcommand{\numlambdaavg}{\ensuremath{4.17\times10^{-9}}}
\newcommand{\numnbavg}{\ensuremath{2.1\times10^{15}}}
\newcommand{\numnbpeak}{\ensuremath{1.0\times10^{16}}}
\newcommand{\numsigmaHtwo}{\ensuremath{1.61\times10^{-20}}}
\newcommand{\numsigmaKr}{\ensuremath{8.96\times10^{-20}}}
\newcommand{\numsigmaratio}{\ensuremath{5.56}}
\newcommand{\numnHtwo}{\ensuremath{3.22\times10^{17}}}
\newcommand{\numnKr}{\ensuremath{3.22\times10^{16}}}
\newcommand{\numtauHtwous}{\ensuremath{80}}
\newcommand{\numtauKrus}{\ensuremath{145}}
\newcommand{\numpKrequal}{\ensuremath{1.8\times10^{-6}}}
\newcommand{\numPionHtwo}{\ensuremath{1.0\times10^{-3}}}
\newcommand{\numPionKr}{\ensuremath{5.8\times10^{-4}}}
\newcommand{\numtransitns}{\ensuremath{83}}
\newcommand{\numTRFns}{\ensuremath{20}}
\newcommand{\numbunchns}{\ensuremath{2.0}}
\newcommand{\numbunchmm}{\ensuremath{4.8}}
\newcommand{\numgapns}{\ensuremath{18}}
\newcommand{\numBfuniform}{\ensuremath{10}}
\newcommand{\numomegapeavg}{\ensuremath{2.6\times10^{9}}}
\newcommand{\numomegapepeak}{\ensuremath{5.7\times10^{9}}}
\newcommand{\numomegapebunch}{\ensuremath{11}}
\newcommand{\numwellavg}{\ensuremath{158}}
\newcommand{\numwellpeak}{\ensuremath{791}}
\newcommand{\numtexpelHtwons}{\ensuremath{44}}
\newcommand{\numtexpelKrns}{\ensuremath{283}}
\newcommand{\numtcrossns}{\ensuremath{1.1}}
\newcommand{\numtescapens}{\ensuremath{107}}
\newcommand{\numetafreeHtwo}{\ensuremath{1\times10^{-3}}}
\newcommand{\numetafreeKr}{\ensuremath{7\times10^{-4}}}
\newcommand{\numionfracHtwo}{\ensuremath{5\times10^{-4}}}
\newcommand{\numionfracKr}{\ensuremath{2\times10^{-3}}}
\newcommand{\numBconfine}{\ensuremath{53}}
\newcommand{\numpicdtfs}{\ensuremath{338}}
\newcommand{\numpicns}{\ensuremath{6.8}}
\newcommand{\numpicdynns}{\ensuremath{41}}
\newcommand{\numpicovertau}{\ensuremath{5\times10^{-4}}}
\newcommand{\numzdoublecm}{\ensuremath{8.6}}
\newcommand{\numzdoublepeakcm}{\ensuremath{3.8}}
\newcommand{\numgrowthKzero}{\ensuremath{5.1}}
\newcommand{\numgrowthKhalf}{\ensuremath{3.4}}
\newcommand{\numgrowthKtenth}{\ensuremath{1.57}}
\newcommand{\numgrowthpeak}{\ensuremath{12.6}}
\newcommand{\numKeffpeakstatic}{\ensuremath{0.82}}
\newcommand{\numEkev}{\ensuremath{30}}
\newcommand{\numIavgma}{\ensuremath{10}}
\newcommand{\numIpeakma}{\ensuremath{50}}
\newcommand{\numabeammm}{\ensuremath{2}}
\newcommand{\numbpipemm}{\ensuremath{10}}
\newcommand{\numLcellcm}{\ensuremath{20}}
\newcommand{\numfRFmhz}{\ensuremath{50}}
\newcommand{\numdphideg}{\ensuremath{36}}
\newcommand{\numBf}{\ensuremath{5}}
\newcommand{\numpHtwo}{\ensuremath{1\times10^{-5}}}
\newcommand{\numpKr}{\ensuremath{1\times10^{-6}}}
\newcommand{\numEeev}{\ensuremath{10}}


% Visible placeholder for content that needs production PIC results.
\newcommand{\pending}[1]{\textcolor{red!70!black}{\textbf{[PENDING:} #1\textbf{]}}}
\newcommand{\Keff}{K_{\mathrm{eff}}}
\newcommand{\etaavg}{\eta_{\mathrm{avg}}}
\newcommand{\Htwo}{\mathrm{H}_2}

\begin{document}

\title{Compact beam-ionized plasma neutralizer for high-current compact-cyclotron
axial injection: design constraints and simulation framework}

\author{Chong Shik Park}
\email{kuphy@korea.ac.kr}
\affiliation{Department of Accelerator Science and Center for Accelerator Research,
Korea University, Sejong, 30019 Republic of Korea}

\date{\today}

\begin{abstract}
High-current compact cyclotrons require multi-mA, low-energy proton beams to be
transported through a short axial injection line whose space-charge forces are
largely uncompensated. We study a compact gas cell, placed between the RF buncher and
the main solenoid, in which the beam ionizes $\Htwo$ or Kr to build a compensating
electron column. For a \numEkev~keV, \numIavgma~mA proton beam of \numabeammm~mm radius
the generalized perveance is $K_0=\numKzero$, and a laminar beam expands by a factor
\numgrowthKzero{} over the \numLcellcm-cm cell length. Using Rudd-model ionization
cross sections we obtain neutralization build-up times of \numtauHtwous~$\mu$s for
$\Htwo$ at $\numpHtwo$~Torr and \numtauKrus~$\mu$s for Kr at $\numpKr$~Torr; Kr reaches
the $\Htwo$ build-up rate at $\numpKrequal$~Torr, not at a tenfold lower pressure.
A time-scale analysis shows that the achievable neutralization is set by electron
confinement rather than by the ionization rate: freely escaping secondary electrons
limit cell-local neutralization to $\eta\sim10^{-3}$, and for a bunched beam
(\numfRFmhz~MHz, bunching factor $B_f$) electrons cross the beam in
\numtcrossns~ns whereas the inter-bunch gap lasts \numgapns~ns, so an axial magnetic
field of order \numBconfine~mT or electrostatic end plugs are required. Even with an
ideal static background, peak-bunch perveance is reduced only to
$\Keff/K_0\approx1-\etaavg/B_f$ (\numKeffpeakstatic{} for $\etaavg=0.9$, $B_f=\numBf$).
We describe a WarpX/PICMI framework, including a custom proton-impact Monte Carlo
collision (MCC) extension, for testing these estimates.
\pending{Self-consistent PIC results for local neutralization, envelope and
inflector transmission.}
\end{abstract}

\maketitle

% =====================================================================
\section{Introduction}
\label{sec:intro}

Compact cyclotrons for isotope production and accelerator-driven applications are
pushed toward average beam currents of several mA and beyond. In a compact machine the
beam is produced by an external ion source and injected axially through a spiral
inflector. At the typical injection energy of tens of keV the beam is strongly
space-charge dominated, and the transverse and longitudinal space-charge forces in the
injection line and on the first turns set practical intensity limits~\cite{Baartman1995}.

In low-energy beam transport (LEBT) lines, positive ion beams ionize the residual gas.
Secondary electrons are trapped in the beam potential, and secondary ions are
expelled, so the beam space charge is progressively compensated~\cite{Gabovich1977,Reiser2008}.
Compensation degrees above 90\% are routinely reported for DC beams in magnetic
LEBTs~\cite{Chauvin2012,BenIsmail2007}, while partial compensation and its dependence
on gas species and pressure have been measured and simulated in detail for both positive
and negative ion beams~\cite{ValerioLizarraga2015,Soloshenko2004}. In cyclotron injector
lines, measured compensation ranges from 0 to 60\%, depending on beam intensity,
radius and pressure~\cite{Winklehner2013}.

Compact axial injection lines differ from these LEBTs in two respects. First, they
are short: there is little drift length over which a compensating electron population
can build up. Second, the beam is bunched upstream of the injection optics, so
the beam potential that traps electrons is modulated at the RF frequency. Dedicated
electron columns, in which the beam ionizes a short, locally pressurized gas section
and the resulting electrons are shaped by external solenoid and electrode fields, have
been proposed and simulated for space-charge compensation in rings~\cite{Freemire2018,Park2020}.

This paper evaluates the same idea for compact-cyclotron injection. We consider a
compact beam-ionized $\Htwo$ or Kr gas cell placed upstream of the main solenoid,
\begin{multline*}
\text{buncher}\rightarrow\text{plasma neutralizer}\rightarrow\text{solenoid}\\
\rightarrow\text{Q1}\rightarrow\text{Q2}\rightarrow\text{spiral inflector},
\end{multline*}
and ask what neutralization it can provide, on what time scale, and under which
confinement conditions. Section~\ref{sec:concept} defines the design point.
Section~\ref{sec:model} develops closed-form estimates for perveance, ionization,
confinement and bunched-beam effects. Section~\ref{sec:methods} describes the
WarpX/PICMI simulation framework and its verification plan. Section~\ref{sec:results}
summarizes the analytical results and the PIC results still required, and
Sec.~\ref{sec:discussion} discusses design implications and limitations.

% =====================================================================
\section{Beamline and neutralizer concept}
\label{sec:concept}

Figure~\ref{fig:layout} shows the baseline layout. The neutralizer is a gas cell of
length $L=\numLcellcm$~cm filled with $\Htwo$ or Kr at a controlled pressure in the
$10^{-6}$--$10^{-5}$~Torr range, with optional local solenoid field and end electrodes
for electron confinement. It is placed after the buncher so that it acts on the beam
before the main solenoid focuses it. Table~\ref{tab:params} lists the design-point
parameters used throughout this paper.

\begin{figure}[tb]
\includegraphics[width=\columnwidth]{fig1_beamline_layout}
\caption{Baseline compact-cyclotron axial injection line with the plasma
neutralizer between the RF buncher and the main solenoid (schematic, not to scale).}
\label{fig:layout}
\end{figure}

\begin{table}[tb]
\caption{Design-point parameters. Derived quantities are computed in
\texttt{plasma\_column.design\_estimates}.}
\label{tab:params}
\begin{ruledtabular}
\begin{tabular}{lll}
Quantity & Symbol & Value \\
\hline
Proton kinetic energy & $E_k$ & \numEkev~keV \\
Velocity & $\beta$, $v$ & \numbeta, $\numvbeam$~m/s \\
Average current & $I_{\mathrm{avg}}$ & \numIavgma~mA \\
Beam edge radius & $a$ & \numabeammm~mm \\
Pipe radius & $b$ & \numbpipemm~mm \\
Generalized perveance & $K_0$ & $\numKzero$ \\
Line charge density & $\lambda$ & $\numlambdaavg$~C/m \\
Beam density & $n_b$ & $\numnbavg$~m$^{-3}$ \\
RF frequency & $f_{\mathrm{RF}}$ & \numfRFmhz~MHz \\
Bunch phase width & $\Delta\phi$ & \numdphideg$^\circ$ \\
Bunching factor & $B_f$ & \numBf \\
Cell length & $L$ & \numLcellcm~cm \\
$\Htwo$ pressure & $p_{\Htwo}$ & $\numpHtwo$~Torr \\
Kr pressure & $p_{\mathrm{Kr}}$ & $\numpKr$~Torr \\
Gas temperature & $T_g$ & 300~K \\
\end{tabular}
\end{ruledtabular}
\end{table}

% =====================================================================
\section{Analytical model}
\label{sec:model}

\subsection{Perveance and envelope growth}

For a round beam of particles with charge $q$, mass $m$ and velocity $v=\beta c$, the
generalized perveance is~\cite{Reiser2008}
\begin{equation}
K_0=\frac{qI}{2\pi\varepsilon_0 m v^3\gamma^3},
\label{eq:K0}
\end{equation}
which gives $K_0=\numKzero$ at the design point. The edge radius $a(z)$ of a uniform
round beam in a field-free drift obeys the envelope equation
\begin{equation}
a''=\frac{\Keff}{a}+\frac{\varepsilon^2}{a^3},
\qquad \Keff = K_0\,(1-\eta_{\mathrm{net}}),
\label{eq:envelope}
\end{equation}
where $\varepsilon$ is the edge emittance and
\begin{equation}
\eta_{\mathrm{net}}=\frac{n_e-n_i}{n_b}
\label{eq:eta}
\end{equation}
is the local net neutralization fraction evaluated in the beam core. For a parallel,
zero-emittance beam, Eq.~\eqref{eq:envelope} integrates exactly. The radius doubles after
\begin{equation}
z_2=\frac{a_0}{\sqrt{2\Keff}}\int_1^2\frac{dx}{\sqrt{\ln x}},
\label{eq:z2}
\end{equation}
which is $\numzdoublecm$~cm for $\Keff=K_0$, i.e.\ shorter than the neutralizer
itself. Figure~\ref{fig:envelope}(a) shows $a(z)$ over the cell length. An
uncompensated beam grows by a factor \numgrowthKzero{}; reducing $\Keff/K_0$ to 0.5 and
0.1 reduces the growth factor to \numgrowthKhalf{} and \numgrowthKtenth{},
respectively. These numbers are upper bounds on the growth, because the laminar
model ignores the emittance term and assumes the full perveance acts from $z=0$, but
they show that useful compensation must be established within a few centimetres of the
cell entrance.

\begin{figure*}[t]
\includegraphics[width=0.85\textwidth]{fig4_envelope_drift}
\caption{Laminar envelope $a(z)/a_0$ of a parallel, zero-emittance beam in a
\numLcellcm-cm drift [Eq.~\eqref{eq:envelope}, $\varepsilon=0$] for several
$\Keff/K_0$. (a)~DC slice at $I_{\mathrm{avg}}$. (b)~Peak slice at
$B_f I_{\mathrm{avg}}$; the dashed line is the static-background bound
$\Keff/K_{0}=1-\etaavg/B_f$ with $\etaavg=0.9$. The slice model neglects the finite
bunch length (\numbunchmm~mm) and overestimates the peak-slice growth.}
\label{fig:envelope}
\end{figure*}

\subsection{Ionization source and cross sections}

Secondary electrons are produced by proton-impact ionization,
$p^+ + X\rightarrow p^+ + X^+ + e^-$, with $X=\Htwo$ or Kr. We use total
ionization cross sections from the Rudd semi-empirical model~\cite{Rudd1985,Rudd1992},
tabulated on a uniform center-of-mass energy grid for use in WarpX MCC
[Fig.~\ref{fig:xs}(a)]. The cross section depends on the projectile velocity. The
tables are therefore evaluated at the 30-keV laboratory energy and stored at the
corresponding center-of-mass energies, $E_{\mathrm{cm}}=E_{\mathrm{lab}}\,m_X/(m_p+m_X)$
(20.0~keV for $\Htwo$ and 29.6~keV for Kr). At 30~keV,
$\sigma_{\Htwo}=\numsigmaHtwo$~m$^2$ and $\sigma_{\mathrm{Kr}}=\numsigmaKr$~m$^2$, a
ratio of \numsigmaratio. At this energy the model lies in the low-energy region where
Rudd \emph{et al.}\ report discrepancies between data sets for some targets~\cite{Rudd1983}.
\pending{Quantitative comparison of the tabulated values with the measurements of
Ref.~\cite{Rudd1983}, and the resulting uncertainty band.}

The ionization rate per proton is $n_g\sigma v$, so the characteristic build-up time
of the electron population is
\begin{equation}
\tau_{\mathrm{ion}}=\frac{1}{n_g\,\sigma\,v}.
\label{eq:tau}
\end{equation}
With $n_g=p/k_BT_g$ this gives $\tau_{\mathrm{ion}}=\numtauHtwous~\mu$s for $\Htwo$ at
$\numpHtwo$~Torr ($n_g=\numnHtwo$~m$^{-3}$) and $\numtauKrus~\mu$s for Kr at
$\numpKr$~Torr [Fig.~\ref{fig:xs}(b)]. Because the cross-section ratio
(\numsigmaratio) is smaller than the pressure ratio (10), Kr at $\numpKr$~Torr builds
up \emph{more slowly} than $\Htwo$ at $\numpHtwo$~Torr. Equal build-up times require
$p_{\mathrm{Kr}}=p_{\Htwo}\sigma_{\Htwo}/\sigma_{\mathrm{Kr}}=\numpKrequal$~Torr.
The single-pass ionization probability across the cell, $n_g\sigma L$, is
$\numPionHtwo$ for $\Htwo$ and $\numPionKr$ for Kr.

\begin{figure*}[t]
\includegraphics[width=0.85\textwidth]{fig2_cross_sections_tau}
\caption{(a)~Rudd-model proton-impact ionization cross sections of $\Htwo$ and Kr
versus proton laboratory energy; dots mark 30~keV. (b)~Ionization build-up time
[Eq.~\eqref{eq:tau}] versus pressure at 300~K. Circle: $\Htwo$ at $\numpHtwo$~Torr;
square: Kr at $\numpKr$~Torr; cross: Kr pressure giving the same $\tau_{\mathrm{ion}}$
as the $\Htwo$ design point.}
\label{fig:xs}
\end{figure*}

\subsection{Neutralization balance and electron confinement}

The commonly used build-up law $\eta(t)=\eta_{\mathrm{ss}}[1-\exp(-t/\tau_{\mathrm{ion}})]$
contains the steady-state value $\eta_{\mathrm{ss}}$ as a free parameter. In a simple
particle balance per unit length,
\begin{equation}
\frac{dN_e}{dt}=\frac{N_p}{\tau_{\mathrm{ion}}}-\frac{N_e}{\tau_c},
\qquad
\eta_{\mathrm{ss}}\simeq\min\!\left(1,\frac{\tau_c}{\tau_{\mathrm{ion}}}\right),
\label{eq:balance}
\end{equation}
$\eta_{\mathrm{ss}}$ is set by the electron confinement time $\tau_c$ in the beam
core. Real equilibria are self-limiting: as $\eta\to1$ the trapping well vanishes, so
$\eta_{\mathrm{ss}}<1$ even for long $\tau_c$~\cite{Gabovich1977}.

Transversely, the uncompensated beam potential confines electrons with a well depth
\begin{equation}
\Delta\Phi=\frac{\lambda}{4\pi\varepsilon_0}\left[1+2\ln\frac{b}{a}\right]
\label{eq:well}
\end{equation}
equal to \numwellavg~V at the average line density, much larger than the
\numEeev-eV secondary-electron energy we assume. Longitudinally, the beam potential
inside a short cell is uniform, so electrons are free to stream out of the gas region.
A \numEeev-eV electron leaves the \numLcellcm-cm cell in about \numtescapens~ns. With
$\tau_c$ equal to this time, Eq.~\eqref{eq:balance} gives a cell-local neutralization of only
$\eta\sim\numetafreeHtwo$ ($\Htwo$) and $\numetafreeKr$ (Kr). The electrons are not lost;
they spread over the whole length between the electron-clearing elements of the line.
The local compensation in a compact cell therefore depends on longitudinal confinement,
for example by electrostatic end plugs, or on the cell being part of a longer
compensated section. This is the central design requirement identified here.

Slow ions are expelled radially by the beam field. An ion born at rest at $r_0=a/2$ in
a uniform beam follows $r(t)=r_0\cosh(t/t_i)$, with
$t_i=\sqrt{2\pi\varepsilon_0a^2M/(e\lambda)}$, and reaches the beam edge after
\numtexpelHtwons~ns ($\Htwo^+$) or \numtexpelKrns~ns (Kr$^+$). The resulting ion
contribution to Eq.~\eqref{eq:eta}, $N_i/N_p\approx t_{\mathrm{exp}}/\tau_{\mathrm{ion}}$,
is $\numionfracHtwo$ for $\Htwo$ and $\numionfracKr$ for Kr. It is negligible while the
beam is far from compensation and grows as the well collapses.

\subsection{RF-bunched beam}

The buncher precedes the neutralizer, so the cell sees bunches with RF period
$T_{\mathrm{RF}}=\numTRFns$~ns. A \numdphideg$^\circ$ phase width corresponds to a
bunch duration of \numbunchns~ns, a bunch length of \numbunchmm~mm, and a gap of
\numgapns~ns. A uniformly filled \numdphideg$^\circ$ bunch would have $B_f=\numBfuniform$.
Here $B_f$ is treated as a free parameter, with $B_f=\numBf$ as the reference value.

If the electron background is static at the average-compensating density $\etaavg n_b$,
the peak-bunch net charge gives
\begin{equation}
\frac{K_{\mathrm{eff,peak}}}{K_{0,\mathrm{peak}}}\simeq1-\frac{\etaavg}{B_f},
\label{eq:peak}
\end{equation}
shown in Fig.~\ref{fig:peak}. For $\etaavg=0.9$ and $B_f=\numBf$ the peak-bunch
perveance is reduced only to \numKeffpeakstatic{} of its uncompensated value
$K_{0,\mathrm{peak}}=\numKzeropeak$. Figure~\ref{fig:envelope}(b) shows that this
reduction barely changes the peak-slice envelope.

\begin{figure}[tb]
\includegraphics[width=0.8\columnwidth]{fig5_peak_perveance}
\caption{Peak-bunch perveance ratio for a static electron background,
Eq.~\eqref{eq:peak}. Dotted lines: $B_f=\numBf$ (reference) and
$B_f=\numBfuniform$ (uniform \numdphideg$^\circ$ bunch).}
\label{fig:peak}
\end{figure}

Equation~\eqref{eq:peak} rests on a static background, which is not guaranteed in either
direction. The electron plasma frequency at the peak bunch density is
$\omega_{pe}=\numomegapepeak$~rad/s, so $\omega_{pe}\Delta t_b\approx\numomegapebunch$.
Electrons can therefore respond within one bunch, and confined electrons could
partially follow the bunch structure and exceed the bound in Eq.~\eqref{eq:peak}. Conversely, during
the \numgapns-ns gap the beam potential is absent while a \numEeev-eV electron crosses
the beam radius in \numtcrossns~ns. Without external confinement, the electron column
accumulated during a bunch is therefore lost before the next bunch arrives, and even
$\etaavg$ cannot be sustained. Holding the electrons transversely during the gap
requires a Larmor radius well below $a$. Taking $r_L=a/10$ for \numEeev-eV electrons gives
$B_z\approx\numBconfine$~mT, which motivates the optional local solenoid field of the
cell.

\subsection{Time-scale hierarchy}

Figure~\ref{fig:timescales} collects the relevant time scales. They span seven orders
of magnitude, from the plasma period to the ionization time. Two consequences
follow. First, the bunch-gap and electron-crossing times make confinement, not
ionization, the limiting process. Second, the ionization time exceeds any affordable
3D PIC run duration by more than three orders of magnitude (Sec.~\ref{sec:methods}),
so the steady state cannot be reached by direct simulation of the build-up.

\begin{figure*}[t]
\includegraphics[width=0.72\textwidth]{fig3_timescales}
\caption{Characteristic time scales at the design point: electron plasma period at
peak beam density, electron crossing and escape times (\numEeev-eV secondaries),
bunch structure, ion expulsion, proton transit, PIC run durations, and ionization
build-up times.}
\label{fig:timescales}
\end{figure*}

% =====================================================================
\section{Simulation methods}
\label{sec:methods}

\subsection{Code and model hierarchy}

Simulations use WarpX~\cite{Vay2021}, built on AMReX~\cite{Zhang2019}, driven through
its PICMI Python interface. Each case is launched as a separate process from a YAML case
file, and a metadata record (repository and WarpX commits, environment, command line,
gas, pressure, geometry, fields, beam and numerical parameters) is written with every
run. Four model levels are distinguished:
\begin{enumerate}
\item \emph{Vacuum reference}: no gas and no compensating species.
\item \emph{Seeded compensation}: an electron/ion population is initialized in the
cell at a density given by the build-up law with a prescribed $\eta_{\mathrm{ss}}$.
This is an analytic source estimate, not a self-consistent ionization model.
\item \emph{Python-callback source}: electron-ion pairs are injected each step at the
analytic rate $n_g\sigma v$ along the beam. This is a data-driven source model.
\item \emph{Custom proton-impact MCC}: a C++ extension of the WarpX background-MCC
module adds an ion-impact ionization process, whose filter and transform functors
create an $e^-$/ion pair from the projectile proton without slowing it.
\end{enumerate}
WarpX's built-in MCC impact ionization is designed for electron-impact
processes~\cite{Vahedi1995}. Level~4 therefore depends on the custom extension, which is
maintained as a patch against WarpX and is not yet validated (Sec.~\ref{sec:verification}).
The present production driver supports levels 1--3, optionally with built-in
electron-impact ionization of the seeded electrons and proton charge exchange on $\Htwo$.

\subsection{Numerical parameters}

The baseline 3D Cartesian domain spans $|x|,|y|\le10$~mm and $-2\le z\le24$~cm on a
$64\times64\times512$ grid, which resolves the beam radius with about six cells. The driver
uses an electromagnetic solver with CFL number 0.5, so $\Delta t=\numpicdtfs$~fs. A
seeded run of $2\times10^4$ steps covers \numpicns~ns, about 0.1 of a cell transit,
and a dynamic-source run of $1.2\times10^5$ steps covers \numpicdynns~ns. The latter is
$\numpicovertau$ of $\tau_{\mathrm{ion}}$ for $\Htwo$. Levels~3 and 4 therefore test source
rates and early-time dynamics, not the steady state.
\pending{Convergence study in grid size, macroparticles per cell and $\Delta t$;
justification of the electromagnetic versus electrostatic solver for this
non-relativistic problem.}

\subsection{Diagnostics}

Global particle counts $N_p$, $N_e$, $N_i$ are recorded every few steps. A global
ratio $N_e/N_p$ includes electrons outside the beam and outside the cell. We
therefore base all compensation claims on the local diagnostic
\begin{equation}
\eta_{\mathrm{local,net}}=\frac{\langle n_e\rangle-\langle n_i\rangle}{\langle n_b\rangle},
\end{equation}
where the averages are taken over $r\le a$ inside the cell. For bunched beams we report
the RF-period average, the peak-bunch value and the local value separately.

\subsection{Verification plan}
\label{sec:verification}

The custom MCC extension will be checked against analytic targets: no ionization
without gas or with $\sigma=0$; a fixed-cross-section rate $dN_e/dt=N_pn_g\sigma v$;
the Kr/$\Htwo$ yield ratio equal to $\sigma_{\mathrm{Kr}}/\sigma_{\Htwo}$ at equal
density; convergence of the per-step probability $P=1-\exp(-n_g\sigma v\Delta t)$ for
$\Delta t$, $\Delta t/2$, $\Delta t/4$; conservation of physical particle number under
macroparticle weighting; and secondary-electron energy bookkeeping.
\pending{Execution of these tests with the patched WarpX build; the WarpX build used
must be shown to match the archived patch.}

% =====================================================================
\section{Results}
\label{sec:results}

\subsection{Analytical design results}

The analytical results of Sec.~\ref{sec:model} are summarized below:
\begin{itemize}
\item An uncompensated \numIavgma-mA beam doubles its radius in \numzdoublecm~cm and
grows by \numgrowthKzero$\times$ over the cell. Compensation must therefore be
established near the cell entrance.
\item Ionization build-up takes \numtauHtwous~$\mu$s ($\Htwo$, $\numpHtwo$~Torr) and
\numtauKrus~$\mu$s (Kr, $\numpKr$~Torr). Kr matches $\Htwo$ at $\numpKrequal$~Torr.
\item Without longitudinal confinement, cell-local neutralization is limited to
$\sim10^{-3}$ by electron escape.
\item For a bunched beam, transverse electron confinement during the \numgapns-ns gap
requires $B_z\sim\numBconfine$~mT or equivalent.
\item With an ideal static background, peak-bunch perveance is reduced only to
\numKeffpeakstatic{} of its uncompensated value ($\etaavg=0.9$, $B_f=\numBf$).
\end{itemize}

\subsection{PIC results}

\pending{The following results require production runs that are not yet available:}
\begin{enumerate}
\item Local $\eta_{\mathrm{local,net}}(t)$ and $\Keff/K_0$ in the cell for $\Htwo$ and
Kr, for the seeded, callback and custom-MCC levels, with and without the cell solenoid
field and end electrodes.
\item Electron confinement time in the cell as a function of $B_z$ and electrode
bias, to test Eq.~\eqref{eq:balance}.
\item Bunched-beam runs resolving the bunch and gap. These would give the RF-averaged and peak-bunch
$\eta$, to be compared with Eq.~\eqref{eq:peak} and with the electron-loss estimate
during the gap.
\item rms size, divergence and emittance at the cell exit, and envelope transport
through the solenoid and Q1/Q2 to the inflector entrance, with transmission into the
inflector acceptance.
\end{enumerate}

% =====================================================================
\section{Discussion}
\label{sec:discussion}

\paragraph{Gas choice.}
Kr offers a larger cross section per atom, but at the pressures usually quoted it does
not build up faster than $\Htwo$. Kr is also heavier, so Kr$^+$ ions stay in the beam
about six times longer, which raises the ion contribution to $\eta_{\mathrm{net}}$
near full compensation. Its gas load on the downstream cyclotron vacuum must also be
assessed. $\Htwo$ is chemically compatible with hydrogen ion sources. The choice
should be made at equal $\tau_{\mathrm{ion}}$ and should include charge-exchange beam
loss, for which measured capture cross sections are available~\cite{Rudd1983} but are
not yet included in our source model.
\pending{Charge-exchange loss and small-angle scattering estimates for both gases.
The Highland formula is not appropriate at keV energies; screened-Coulomb elastic
cross sections are required.}

\paragraph{Placement.}
Placing the cell before the solenoid lets it act on the beam before magnetic focusing.
The solenoid fringe field can contribute to transverse electron confinement, whereas
electrostatic elements and the buncher gap clear electrons and bound the compensated
section. Alternative placements are outside the baseline design and are not considered
here.

\paragraph{Overcompensation.}
Values of $\Keff/K_0<0$ correspond to net negative charge in the core and focus the
beam. They can arise transiently in seeded models with $\eta>1$ or in strongly confined
columns. Such values are always labeled explicitly.

\paragraph{Limitations.}
All quantitative statements in this draft are closed-form estimates for a uniform
round beam. The envelope model is laminar and slice-based. The confinement estimate
uses a single secondary-electron energy of \numEeev~eV. The bunched-beam bound assumes
a static background. Cross sections come from a semi-empirical model at the edge of its
validated range. These estimates define what the PIC campaign must test; they do not
replace it.

% =====================================================================
\section{Conclusion}
\label{sec:conclusion}

A compact beam-ionized gas cell upstream of the main solenoid is a plausible way to
reduce the perveance of a multi-mA, 30-keV proton beam in a compact-cyclotron injection
line. However, its performance is controlled by electron confinement rather than by
the ionization rate or the gas species. For a DC beam, longitudinal confinement is needed
to raise cell-local neutralization above the free-escape level. For the RF-bunched beam
that actually enters the cell, transverse confinement during the inter-bunch gap is
required, and even ideal average compensation leaves most of the peak-bunch space charge.
The WarpX/PICMI framework described here, with a custom proton-impact MCC extension,
is designed to quantify these effects.
\pending{Final conclusions after the PIC campaign.}

\begin{acknowledgments}
\pending{Funding acknowledgments.}
\end{acknowledgments}

\appendix
\section{Reproducibility}
All numbers and figures in this paper are generated by
\texttt{python scripts/make\_manuscript\_figures.py} from the module
\texttt{plasma\_column.design\_estimates} and the tabulated cross sections in the
project repository. The script records the repository commit and environment in
\texttt{paper/manuscript/manuscript\_numbers.json}.

\bibliography{references}

\end{document}
